\exercisesfor{4}

\begin{enumerate}
\def\labelenumi{\arabic{enumi}.}
\tightlist
\item
  设 G 是一个李群, H 是 G 的正规李拟子群, \(\pi\) 是从 G 到 G/H 的典范映射. 在 G/H 上存在唯一的李群结构, 具有如下性质: 对于从 G/H 到李群 G' 的同态 \(\theta\), \(G'\) 是李群态射当且仅当 \(\theta \circ \pi\) 是李群态射. 此外, L(G/H) 典范地同构于 L(G)/L(H). (令 Q 为满足 L(Q) = L(G)/L(H) 的李群芽. 证明必要时缩小 Q 后, Q 可识别为 G/H 中 0 的一个开邻域, 从而 §1 的命题 18 可应用于 G/H.)
\end{enumerate}

¶ 2. 设 G 和 H 是两个李群, f 是从 G 到 H 的李群态射.

\begin{enumerate}
\def\labelenumi{(\roman{enumi})}
\item
   的核 \(\mathbf{N}\) 是 \(f\)\(\mathbf{G}\) 的正规李拟子群, 且 \(\mathbf{L}(\mathbf{N}) = \operatorname{Ker} \mathbf{L}(f)\) . (使用 \(\mathbf{G}\) 和 \(\mathbf{H}\) 的指数映射.)
\item
  令 g \(g \colon \mathrm{G} / \mathrm{N} \to f(\mathrm{G})\) 为 f 取商后导出的映射. 若 f 和 \(f\)\(f\)\(\mathrm{L}(f)\) 的像都是闭的, 且 \(\mathrm{G}\) 的拓扑有可数基, 则 f(\(f(\mathrm{G})\)G) 是 \(\mathrm{H}\) 的李拟子群, 其李代数为 \(\operatorname{Im} \mathrm{L}(f)\), 且 g 是李群同构, 其中 \(g\)\(\mathrm{G} / \mathrm{N}\) 具有练习 1 中定义的结构. (利用 (i) 将问题归约为 \(\mathrm{N} = \{e\}\) 的情形.)
\end{enumerate}

\begin{enumerate}
\def\labelenumi{\arabic{enumi}.}
\setcounter{enumi}{2}
\tightlist
\item
  设 G 是一个李群, U 是 L(G) 中 0 的开邻域, \(\phi\) 是从 U 到 G 的解析映射, 满足 \(\phi(0) = e\) 且 \(T_{0}\phi = \mathrm{Id}_{\mathrm{L}(G)}\). 以下条件等价:
\end{enumerate}

\begin{enumerate}
\def\labelenumi{(\alph{enumi})}
\item
  \(\phi\) 是指数映射;
\item
  存在不同于 0, 1, -1 的 \(m \in \mathbf{Z}\), 使得在 0 的某个邻域内 \(\phi(mx) = \phi(x)^m\).
\end{enumerate}

(若条件 (b) 成立, 证明 \(\phi^{*}\) 满足定理 4 的条件 (iii).)

¶ 4. (a) 设 K = R 或 C, 或剩余域特征为 \textgreater0 的超度量域. 设 G 是 K 上的李群, U 是 L(G) 中 0 的开邻域, \(\phi: U \to G\) 是在 0 处可微的映射, 并且
满足 \(\phi(0) = e\), \(T_0(\phi) = \mathrm{Id}_{L(G)}\), 且当 \(\phi((\lambda + \lambda')b) = \phi(\lambda b)\phi(\lambda'b)\)\(\lambda b\), \(\lambda'b\), \((\lambda + \lambda')b\) 属于 U 时有 . 则 \(\phi\) 在 0 的某个邻域上与一个指数映射重合.

\begin{enumerate}
\def\labelenumi{(\alph{enumi})}
\setcounter{enumi}{1}
\tightlist
\item
  设 k 是特征为 0 的域. 令 \(k\)\(\mathbf{K}\) 为赋值域 \(k((\mathbf{X}))\), 它的特征为 0. 令 \(\mathbf{G}\) 为加法李群 \(k[[\mathbf{X}]]\). 设 \(\phi\) 是从 \(k\)\(\mathbf{G}\) 到 \(\mathbf{G}\) 的连续 k-线性映射, 满足对每个整数  都有 \(\phi(\mathbf{X}^n) = \mathbf{X}^n + \mathbf{X}^{2n}\). 则 \(n \geqslant 0\)\(\phi\) 满足 (a) 的条件, 但在 0 的任何邻域上都不与指数映射重合.
\end{enumerate}

\begin{enumerate}
\def\labelenumi{\arabic{enumi}.}
\setcounter{enumi}{4}
\tightlist
\item
  设 \((e_1, e_2, e_3)\) 是 \(\mathbf{R}^3\) 的典范基. 在 \(\mathbf{R}^3\) 上赋予幂零李代数结构 \(\mathfrak{g}\), 使得 \([e_1, e_2] = e_3\), \([e_1, e_3] = [e_2, e_3] = 0\). 令 \(c_{ijk}\) 表示相对于基  的 \(\mathfrak{g}\) 的结构常数.\((e_1, e_2, e_3)\)
\end{enumerate}

\begin{enumerate}
\def\labelenumi{(\alph{enumi})}
\tightlist
\item
  证明在 \(\mathbf{R}^3\) 上微分形式
\end{enumerate}

\[
\omega_ {1} = d x _ {1}, \quad \omega_ {2} = d x _ {2}, \quad \omega_ {3} = - x _ {1} d x _ {2} + d x _ {3}
\]

满足关系

\[
d \omega_ {k} = - \sum_ {i <   j} c _ {i j k} \omega_ {i} \wedge \omega_ {j} \quad (k = 1, 2, 3)
\]

并且在 \(\mathbf{R}^3\) 的每一点线性无关.

\begin{enumerate}
\def\labelenumi{(\alph{enumi})}
\setcounter{enumi}{1}
\tightlist
\item
  推出在 \(\mathbf{R}^3\) 中 0 的某个开邻域上存在李群芽结构 \(\mathbf{G}\), 使得 \(\mathrm{L}(\mathrm{G}) = \mathfrak{g}\), 且
\end{enumerate}

\[
(a _ {1}, a _ {2}, a _ {3}) (x _ {1}, x _ {2}, x _ {3}) = (x _ {1} + a _ {1}, x _ {2} + a _ {2}, x _ {3} + a _ {3} + a _ {1} x _ {2})\tag{1}
\]

当 \((a_{1}, a_{2}, a_{3})\), \((x_{1}, x_{2}, x_{3})\) 充分接近 0 时成立. 证明事实上公式 (1) 在 \(\mathbf{R}^{3}\) 上定义了幂零李群结构.

¶ 6. 设 G 是有限维李群, g 是其李代数, g* 是 g 的对偶向量空间, σ 是 G 在 g 上的伴随表示, ρ 是 G 在 g* 上的表示, 对所有 g ∈ G 定义为 ρ(g) = tσ(g)-1.

\begin{enumerate}
\def\labelenumi{(\alph{enumi})}
\item
  对所有 , \(\mathrm{L}(\rho)(x) = -^{t}(\mathrm{ad} x)\).\(x \in g\)
\item
  设 \(f \in \mathfrak{g}^*\). 令 \(G_f\) 表示 f 在 G 中的稳定子; 它是 G 的李子群, 且 \(f\)\(G\)\(G\)\(\mathfrak{g}_f = L(G_f)\) 是满足  的 \(x \in G\) 的集合.\({}^t (\operatorname{ad}x)f = 0\)
\item
  对 \(x, y\), 记 \(\mathfrak{g}\)\(\mathrm{B}_f(x, y) = f([x, y])\). 证明 \(\mathrm{B}_f\) 是 \(\mathfrak{g}\) 上的交错双线性形式, 与 \(s_f\) 左相关的从 \(\mathfrak{g}\) 到 \(\mathfrak{g}^*\) 的映射  是映射  , 且相对于 \(\mathrm{B}_f\), \(x \mapsto {}^t (\mathrm{ad}x)\)\(f\)\(\mathfrak{g}\) 的正交补是 \(\mathrm{B}_f\)\(\mathfrak{g}_f\). 令 \(\beta_f\) 表示在取商时由  导出的 \(\mathfrak{g} / \mathfrak{g}_f\) 上的非退化交错双线性形式 (因而其维数为偶数). 我们将  记为  上 \(\mathrm{B}_f\) 的逆形式.\(\beta_f'\)\(\mathrm{B}_f\)\((\mathfrak{g} / \mathfrak{g}_f)*\)
\item
  令 \(\Omega\) 表示 G 在 \(G\)\(\mathfrak{g}^*\) 上的一个轨道; 通过结构传递, 赋予它由 \(G / G_f\) 上的结构导出的流形结构, 其中 f 是 \(f\)\(\Omega\) 的任意一点. 然后 G 在 \(G\)\(\Omega\) 上解析地左作用. 令 D 为相应的无穷小作用法则. 对所有 \(D\)\(a \in \mathfrak{g}\) 和所有 \(f \in \Omega\),
\end{enumerate}

\[
\mathrm{D} _ {a} (f) = - (^ {t} \mathrm{ad} a) f.
\]

切子空间 \(\mathrm{T}_{f}(\Omega)\) 是 \(s_{f}\) 的像, 即  中 \(g_{f}\) 的正交补, 并且典范地与 \(g^{*}\)\((\mathfrak{g}/\mathfrak{g}_{f})^{*}\) 对应, 因而 \(s_{f}\) 定义了从  到其对偶的典范同构 \(r_{f}\). 场 \(\mathrm{T}_{f}(\Omega)\)\(f \mapsto \beta_{f}'\) 是  上的 2 次解析微分形式 \(\omega\), 且在 G 下不变. 对于 \(\Omega\)\(f \in \Omega\), \(a \in g\), \(b \in g\), \(\omega(\mathrm{D}_{a}(f), \mathrm{D}_{b}(f)) = f([a, b])\).

\begin{enumerate}
\def\labelenumi{(\alph{enumi})}
\setcounter{enumi}{4}
\item
  证明 \(d\omega = 0\) .
\item
  若 \(\alpha\) 是 \(\Omega\) 上的 1 次微分形式, 同构 \(r_f\) 使我们能够将 \(\alpha\) 识别为一个向量场. 令 A 为取值于 K 的 \(\Omega\) 上解析函数的集合. 对于 A 中的 \(\phi, \psi\), 记 \([\phi, \psi] = \omega(d\phi, d\psi) \in A\) . 证明如此定义的 A 具有李代数结构.
\item
  对所有 \(a \in \mathfrak{g}\), 令 \(\psi_{a}\) 为  上的函数 \(f \mapsto f(a)\). 证明 \(\Omega\)\(a \mapsto \psi_{a}\) 是从 \(\mathfrak{g}\) 到李代数 A 的同态, 且形式 \(d\psi_{a}\) 通过同构 \(r_{f}\) 与向量场 \(\mathrm{D}_{a}\) 对应.
\item
  令 U 是 \(\mathfrak{g}^*\) 的开子集, \(\phi\) 是 U 上的解析函数. 对于所有 \(f\in \mathrm{U}\), \(d_{f}\phi\)\(\phi\) 在 \(f\) 处的微分  被识别为 \(\mathfrak{g}\) 的一个元素. 假设对 G 在  上的作用所定义的每个向量场 X 都有 \(\mathrm{X}\phi = 0\). 证明于是对所有 \(\mathfrak{g}^*\)\(f\in \mathrm{U}\), \(d_{f}\phi\) 属于 \(\mathfrak{g}_f\) 的中心.
\item
  对所有 \(f \in \mathfrak{g}^*\), 记 \(r_f = \dim \mathfrak{g}_f\). 令 \(r = \inf_{f \in \mathfrak{g}^*} r_f\). 证明集合 V = \(V\) 在 \(f \in \mathfrak{g}^*\)\(r_f = r\)\(\mathfrak{g}^*\) 中开且稠密. 证明对所有 \(f \in V\), \(\mathfrak{g}_f\) 是交换的. (在 f 的某个邻域内构造 r 个满足 (h) 条件的函数 \(r\)\(\phi\), 并使它们在 f 处的微分线性无关.)\((h)\)\(f\)\(f\)
\end{enumerate}

¶ 7. (a) 设 \((\lambda, x) \mapsto \lambda.x\) 是 R 在 T 上的连续左作用法则. 则以下两种情形之一成立:

\begin{enumerate}
\def\labelenumi{(\arabic{enumi})}
\item
  要么存在一个不动点, 且每一点的稳定子或者是 \(\mathbf{R}\), 或者是 \(\{0\}\) ;
\item
  要么存在一点, 其稳定子是 \(\mathbf{R}\) 的无限离散子群, 且 \(\mathbf{R}\) 在 \(\mathbf{T}\) 上传递作用. (若 \(\{\lambda \in \mathbf{R}|\lambda .x = x\} = \{0\}\), 则 x 的轨道同胚于 \(x\)\(\mathbf{R}\), 并具有边界点 \(\lim_{\lambda \to +\infty}\lambda .x\), 这些边界点是不动点. 若 \(\{\lambda \in \mathbf{R}|\lambda .x = x\} = \mathbf{Z}a\) 且 \(a\neq 0\), 则 x 的轨道同胚于 \(x\)\(\mathbf{T}\), 且 \(\mathbf{R}\) 是传递的.)
\end{enumerate}

\begin{enumerate}
\def\labelenumi{(\alph{enumi})}
\setcounter{enumi}{1}
\item
  若 f 是 \(f\)\(\mathbf{T}\) 到自身的同胚, k 是一个 \(k\)\(\geqslant 1\) 的整数, 则称 \(x \in \mathbf{T}\) 是 f 的周期为 k 的周期点, 如果 \(k\)\(f\)\(f^k(x) = x\) 且对  有 \(f^h(x) \neq x\). 在 (a) 的记号下, 令 \(1 \leqslant h < k\)\(f_{\lambda}\) 为映射 \(x \mapsto \lambda .x\) 所定义的同胚. 则 \(k > 1\)\(f_{\lambda}\) 的周期大于 1 且等于 k 的周期点集为空或等于 \(\mathbf{T}\). (若  存在周期为 \(k > 1\) 的周期点, 则其稳定子不同于 \(f_{\lambda}\)\{0\} 和 \(\mathbf{R}\), 因而 \(\mathbf{R}\) 在 \(\mathbf{T}\) 上传递作用, 且 \(\mathbf{T}\) 的所有点都是 \(k\) 的周期为 k 的周期点.)\(f_{\lambda}.\)
\item
  令 \(\Omega\) 是 \(e\)\(\mathrm{Diff}^{\infty}(\mathbf{T})\) 中 e 的一个邻域 (《可微与解析流形》, R, 15.3.8, 注 (1)). 存在整数 \(k > 1\) 以及无不动点的 \(f\in \Omega\), 使得 f 的周期为 k 的周期点集既非空又不等于 \(k\)\(f\)\(\mathbf{T}\). (令 \(p\colon \mathbf{R}\to \mathbf{T}\) 为典范映射. 令 \(\rho\) 为  上的映射 \(p(x) \mapsto p\left(x + \frac{2\pi}{k}\right)\). 令 \(\mathbf{T}\)\(\mathbf{T}\)\(\Omega'\) 是 \(e\)\(\text{Diff}^\infty(\mathbf{T})\) 中 e 的一个邻域, 使 \(\Omega'^2 \subset \Omega\). 当 k 充分大时, \(k, \rho \in \Omega'\). 另一方面, 令 \(\sigma\) 是 \(\mathbf{T}\) 到自身的映射, 使得 \(\mathbf{T}\)\(\sigma(y) = y\) 对于
\end{enumerate}

\[
y \notin p \left(\right] 0, \frac {2 \pi}{k} [)
\]

且当  时 \(\sigma(p(x)) = p(x + \lambda(x))\), 其中\(x \in \left]0, \frac{2\pi}{k}\right\}\)

\[
0 <   \lambda (x) <   \frac {2 \pi}{k} - x;
\]

可以选取 \(\sigma \in \Omega'\). 则 \(f = \rho \circ \sigma\) 具有所要求的性质.) 根据 (b), 这样的 f 不可能属于从 \(f\)\(\mathbf{R}\) 到 \(\operatorname{Diff}^{\infty}(\mathbf{T})\) 的连续态射的像.

\begin{enumerate}
\def\labelenumi{(\alph{enumi})}
\setcounter{enumi}{3}
\item
  设 Y 是维数 \(\geqslant1\) 的紧可微流形. \(\operatorname{Diff}^{\infty}(Y)\) 的每个 e 邻域都包含一个具有如下性质的元素 h: h 不属于从 R 到 \(\operatorname{Diff}^{0}(Y)\) 的任何连续态射的像. (流形 Y 包含一个形如 \(T \times D\) 的开子流形, 其中 D 是 \(R^{n}\) 中半径为 1 且中心为 0 的开欧氏球 ( \(n = \dim Y - 1\) ). 取 \(g \in \operatorname{Diff}^{\infty}(D)\), 使得 \(g(0) = 0\), 当  时 \(\|g(x)\| > \|x\|\), 当 \(0 < \|x\| < \frac{1}{2}\) 时 \(g(x) = x\), 且 g 在 \(\|x\| \geqslant \frac{1}{2}\)\(\operatorname{Diff}^{\infty}(D)\) 中充分接近 e. 令 f 如 (e) 中, 对 \textless, 令 {}\textless{}\(f_{i}: T \to T\) 如下定义: 对于 \(x \in T\), \(y \in p^{-1}(x)\), \(z \in p^{-1}(f(x))\) 且 \(|z - y| < \pi\), 记 \(f_{i}(x) = p(ty + (1 - t)z)\). 则存在 \(h \in \operatorname{Diff}^{\infty}(Y)\), 使得对 \(h(x, y) = (f_{2\|x\|}(x), g(y))\)\(x \in T\), \(y \in D\) 有 , 且对  有 \(h(u) = u\); 此外, 可以选择 f 和 g, 使 h 在 \(u \in Y - (T \times D)\)\(\operatorname{Diff}^{\infty}(Y)\) 中任意接近 e. \(T \times \{0\}\) 中的点正是这样的 \(y \in Y\): 当 n 趋于  时, \(h^{kn}(x)\) 趋于一个不被 h 固定的点 (事实上它是周期为 k 的周期点). 因此, 每个与 h 交换的 Y 到自身的同胚都保持 \(T \times \{0\}\) 不变. 然后应用 (e).
\item
  设 Y 是维数 \(\geqslant1\) 的紧可微流形. 不存在李群 G 以及连续态射 \(\operatorname{Diff}^{\infty}(Y)\to G\), \(G\to\operatorname{Diff}^{0}(Y)\), 使它们的复合为 \(\operatorname{Diff}^{\infty}(Y)\) 到 \(\operatorname{Diff}^{0}(Y)\) 的典范嵌入. (使用 (d).)
\end{enumerate}

\begin{enumerate}
\def\labelenumi{\arabic{enumi}.}
\setcounter{enumi}{7}
\tightlist
\item
  设 G 是有限维李群. 假设 L(G) 是单的. 设 A 是 G 的正规子群. 若 A 不是开的, 则 A 是离散的, 且其在 G 中的中心化子是开的. (令 a 为 A 在 e 处相切的子代数. 证明 \(a = \{0\}\). 取 A 中不属于 G 中心的元素 x. 考虑从 G 到 A 的映射 \(y \mapsto xyj^{-1}\). 由 \(a = \{0\}\) 推出 G 中 x 的中心化子是开的. 然后利用指数映射证明, x 在 G 中的中心化子包含一个与 x 无关的 e 邻域.)
\end{enumerate}

¶ 9. 设 G 是 K 上维数 n 的紧李群.~假设李代数 L(G) 是单的. 对所有 \(g \in G\), 每个整数 \(m \geqslant 1\) 以及 G 中 e 的每个邻域 V, 令 M(g, m, V) 表示 G 中形如 \(\prod_{i=1} x_i(y_i, g) x_i^{-1}\) 的元素的集合, 其中 \(x_1, \ldots, x_m, y_1, \ldots, y_m\) 属于 V.

\begin{enumerate}
\def\labelenumi{(\alph{enumi})}
\tightlist
\item
  证明若 g 是中心化子不开的 G 中元素, m 是一个 \(g\)\(G\)\(m\)\(\geqslant n\) 的整数且 V 是 e 的邻域, 则 M(g, m, V) 是 e 的邻域. (存在 \(V\)\(e\)\(M(g, m, V)\)\(e\)\(a \in L(G)\), 使 \((\operatorname{Ad} g)(a) \neq a\). 令 \(\phi\) 为 G 的一个指数映射. 令 y 为 0 处映射 \(G\)\(y\)\(\lambda \mapsto \langle \phi(xa), g \rangle\) 的切映射对 1 的像 (其中 \(\lambda \in K\)). 则 \(b = (\operatorname{Ad} g - 1)(a) \neq 0\). 存在 V 中的 \(x_1, \ldots, x_m\), 使 \(V\) 包含 L(G) 的一个基. 考虑映射\(\{(\operatorname{Ad} x_1)(b), \ldots, (\operatorname{Ad} x_m)(b)\}\)\(L(G)\)
\end{enumerate}

\[
\left(x _ {1} ^ {\prime}, \dots , x _ {m} ^ {\prime}, y _ {1}, \dots , y _ {m}\right) \mapsto \prod_ {i = 1} ^ {n} x _ {i} ^ {\prime} (y _ {i}, g) x _ {i} ^ {\prime - 1}
\]

从 \(\mathbf{G}^{2m}\) 到 G. 证明在 \((x_{1},\ldots ,x_{m},e,\ldots ,e)\) 处的切映射是满射.)

\begin{enumerate}
\def\labelenumi{(\alph{enumi})}
\setcounter{enumi}{1}
\item
  令 \(\mathbf{U}\) 为 \(e\)\(\mathbf{G}\) 中 e 的一个邻域. 令 m 为一个 \(m\)\(\geqslant 1\) 的整数. 证明存在一个中心化子不开的 \(g\)\(\mathbf{G}\) 中元素 g, 使得 \(\mathbf{M}(g, m, \mathbf{G}) \subset \mathbf{U}\). (利用 \(\mathbf{G}\) 的紧性作反证.)
\item
  设 \(\mathbf{G}'\) 是其李代数为单的 K 上的紧李群. 若 \(\mathbf{K}\)\(\phi: \mathbf{G} \to \mathbf{G}'\) 是抽象群的双射同态, 则 \(\phi\) 连续. (对 \(\mathbf{G}'\) 应用 (b), 对 \(\mathbf{G}\) 应用 (a).)
\end{enumerate}

\begin{enumerate}
\def\labelenumi{\arabic{enumi}.}
\setcounter{enumi}{9}
\item
  设 G 是 K 上的李群. 证明存在 e 的一个邻域 V, 具有如下性质: 对 V 中元素的任意序列 \(G\)\(K\)\(V\)\(e\)\((x_{n})\), 若归纳定义 \(V\)\(y_{1} = x_{1}, y_{n} = (x_{n}, y_{n-1})\), 则序列 \((y_{n})\) 趋于 e.\(e\)
\item
  \begin{enumerate}
  \def\labelenumii{(\alph{enumii})}
  \tightlist
  \item
    对于  中的 \(x, y\), 定义 \(\mathbf{S}_{2n-1} \subset \mathbf{C}^n\)\(\alpha(x, y) \in [0, \pi]\) 为
  \end{enumerate}
\end{enumerate}

\[
\cos \alpha (x, y) = \mathscr {R} ((x | y)).
\]

对于  中的 \(s, t\), 记 . 证明 d 是 \(\mathbf{U}(n)\) 上左右不变的距离.\(d(s, t) = \sup_{x \in \mathbf{S}_{2n-1}} \alpha(sx, tx)\)\(d\)\(\mathbf{U}(n)\)

\begin{enumerate}
\def\labelenumi{(\alph{enumi})}
\setcounter{enumi}{1}
\item
  设 \(s \in \mathbf{U}(n)\). 令 \(\theta_1, \ldots, \theta_m\) 为 \(\mathbf{J} - \pi, \pi \mathbf{j}\) 中的数, 使 \(e^{i\theta_j}\) 是 s 的互异特征值. 令 \(s\)\(\theta(s) = \sup_{1 \leqslant j \leqslant m} |\theta_j|\). 则 \(d(e, s) = \theta(s)\). (令 \(V_j\) 为 s 对应于 \(s\)\(e^{i\theta_j}\) 的特征子空间. 对 \(x \in S_{2n-1}\), 通过按  分解, 求 \(\mathcal{R}((x|sx))\) 的一个下界.)\(V_j\)
\item
  设 \(s, t\), 满足 \(\mathbf{U}(n)\)\(\theta(t) < \frac{\pi}{2}\), 且 s 与 \(s\)\((s, t)\) 交换. 则 s 与 t 交换. (在 (b) 的记号下, 令 \(s\)\(t\)\(\mathrm{V}_j'\) 为 \(\mathrm{V}_j\) 的正交补, \(\mathrm{W}_j = t(\mathrm{V}_j)\); 证明 \(\mathrm{W}_j = (\mathrm{W}_j \cap \mathrm{V}_j) + (\mathrm{W}_j \cap \mathrm{V}_j')\), 然后证明 \(\mathrm{W}_j \cap \mathrm{V}_j' = \{0\}\).)
\item
  证明存在  中 e 的一个紧邻域 \(\mathbf{V}\), 它具有练习 10 的性质, 在 \(e\)\(\mathbf{U}(n)\) 的内自同构下不变, 并且若 x, y 是 \(\mathbf{U}(n)\)\(x, y\)\(\mathbf{V}\) 中不可交换的元素, 则 x 与 \(x\)\((x, y)\) 不可交换. (使用 (c).)
\item
  令 \(\beta\) 为 \(\mathbf{U}(n)\) 的归一化 Haar 测度. 令 \(\mathrm{V}_1\) 是 \(e\)\(\mathbf{U}(n)\) 中 e 的对称紧邻域, 使 \(\mathrm{V}_1^2 \in \mathrm{V}\). 令 p 为满足 \(p\) 的整数. 证明对每个 \(\beta (\mathrm{V}_1) > 1 / p\)\(\mathbf{GL}(n,\mathbf{C})\) 的有限子群 F, 都存在 F 的交换正规子群 A, 使 \(\operatorname{Card}(\mathrm{F} / \mathrm{A}) \leqslant p\) (Jordan 定理). (将其归约为 \(\mathrm{F} \subset \mathbf{U}(n)\) 的情形. 取 A 为由 \(\mathrm{F} \cap \mathrm{V}\) 生成的 F 的子群, 并使用 (d).)
\end{enumerate}

\begin{enumerate}
\def\labelenumi{\arabic{enumi}.}
\setcounter{enumi}{11}
\tightlist
\item
  设 G 是有限维连通实或复李群, \(\alpha\) 是 G 上 p 次左不变微分形式. \(\alpha\) 同时右不变, 当且仅当 \(d\alpha = 0\). (注意右不变性的条件可写成
\end{enumerate}

\[
\wedge^ {p} (^ {t} \mathrm{Ad} (s)). \alpha (e) = \alpha (e)
\]

对所有 \(s \in G\) 都成立, 因而等价于条件

\[
\sum_ {j = 1} ^ {f} \left\langle \alpha (e), u _ {1} \wedge \dots \wedge u _ {j - 1} \wedge [ u, u _ {j} ] \wedge u _ {j + 1} \wedge \dots \wedge u _ {p} \right\rangle = 0
\]

对 \(u, u_1, \ldots, u_p\)\(\mathrm{L}(\mathrm{G})\) 中的所有  都成立. 特别地, \(\mathbf{G}\) 上 1 次左右不变微分形式构成维数为

\[
\dim \mathrm{L} (\mathrm{G}) - \dim [ \mathrm{L} (\mathrm{G}), \mathrm{L} (\mathrm{G}) ].
\]

\begin{enumerate}
\def\labelenumi{\arabic{enumi}.}
\setcounter{enumi}{12}
\tightlist
\item
  令 \(\mathbf{R}^n\) 带有通常的内积 \(((\xi_i), (\eta_i)) \mapsto \sum_{i=1} \xi_i \eta_i\). 令 \(\mathbf{I}(n)\) 为 \(\mathbf{R}^n\) 的等距仿射变换群. 它典范地同构于 \(\mathbf{GL}(n + 1, \mathbf{R})\) 的李子群, 该子群由矩阵
\end{enumerate}

\[
S = \left( \begin{array}{c c} U & x \\ 0 & 1 \end{array} \right)
\]

其中 \(U \in \mathbf{O}(n)\), x 是 \(x\)\(\mathbf{R}^n\) 的任意元素 (类型为 \((n,1)\) 的矩阵). \(\mathbf{I}(n)\) 可识别为由 \(\mathbf{O}(n)\) 到  的典范嵌入所定义的  与 \(\mathbf{R}^n\) 的半直积 (此时矩阵 S 记作 \(\mathbf{O}(n)\)\(\mathbf{GL}(n,\mathbf{R})\)\(S\)\((U,x)\)). 令 \(p\colon \mathbf{I}(n)\to \mathbf{O}(n)\) 为典范满射, 核为 \(\mathbf{R}^n\).

\begin{enumerate}
\def\labelenumi{(\alph{enumi})}
\tightlist
\item
  设 \(\Gamma\) 是 \(\mathrm{I}(n)\) 的离散子群. 考虑 \(\tilde{p}(\Gamma)\)\(p(\Gamma)\) 在 \(\mathbf{O}(n)\) 中的闭包 \(V\). 证明其单位元连通分支是交换的. (令 V 为 \(I\)\(\mathbf{O}(n)\) 中 I 的一个邻域, 具有练习 11 (d) 的性质, 并且还满足对所有 \(d\) 有 \(\|U - I\| \leqslant \frac{1}{4}\) (\(U \in V\)\(\operatorname{End}(\mathbf{R}^n)\) 上的范数由 \(\mathbf{R}^n\) 上的欧氏范数导出). 反证, 假设存在 \(S_1 = (U_1, x_1)\), \(S_2 = (U_2, x_2)\), 其中 \(U_1\) 和 \(U_2\) 属于 V 且不交换. 证明 \(V\)\((S_1, S_2) = ((U_1, U_2), y)\), 其中
\end{enumerate}

\[
\| y \| \leqslant \frac {1}{4} (\| x _ {1} \| + \| x _ {2} \|).
\]

\(S_{k}\). 然后对  归纳定义 \(S_{k}=(S_{1},S_{k-1})\). 利用 V 的选取, 证明此序列有无穷多个不同的项, 且在 \(k\geqslant3\) 中有界, 矛盾.)\(\mathbf{I}(n)\)

\begin{enumerate}
\def\labelenumi{(\alph{enumi})}
\setcounter{enumi}{1}
\item
  若 \(\Gamma\) 是晶体学群, 即 \(\mathrm{I}(n)/\Gamma\) 是紧的, 则证明对所有 \(x \in \mathbf{R}^n\), 由  生成的  的仿射子空间 \(\mathbf{L}\) 等于 \(\mathbf{R}^n\)\(\Gamma x\)\(\mathbf{R}^n\). (否则, 对所有 \(y \in \mathbf{R}^n\), \(\Gamma y\) 的所有点到 \(\mathbf{L}\) 的距离都相同; 由于这个距离可以任意大, \(\mathrm{I}(n)/\Gamma\) 不是紧的.)
\item
  若 \(\Gamma\) 是交换的晶体学群, 则 \(\Gamma \subset \mathbf{R}^n\). (若 \(S = (U, x) \in \Gamma\), 使得 \(\mathrm{V} \subset \mathbf{R}^n\) 中由 U 的不动点组成的向量子空间 \(U\) 不等于 \(\mathbf{R}^n\), 且 \(\mathrm{V}^\perp\) 表示 \(\mathrm{V}\) 的正交子空间, 则对所有 \(S' = (U', x') \in \Gamma\), U' 都保持 \(U'\)\(\mathrm{V}\) 和 \(\mathrm{V}^\perp\) 不变. 另一方面, 由于 \(U|\mathrm{V}^\perp\) 没有非零不动点, 证明改变原点后可设 \(x \in \mathrm{V}\). 根据 (b), 存在 \(S' = (U', x') \in \Gamma\), 使 x' 在 \(y'\)\(x'\)\(\mathrm{V}^\perp\) 上的正交投影 \(\neq 0\). 用公式 \(Sx'\)\(S = S'SS'^{-1}\) 计算 Sx', 推出 \(U.y' = y'\), 这与 \(\mathrm{V}\) 的定义矛盾.
\item
  若 \(\Gamma\) 是晶体学群, 则 \(\Gamma \cap \mathbf{R}^n\) 是秩为 n 的自由交换群, 且 \(n\)\(\Gamma / (\Gamma \cap \mathbf{R}^m)\) 是有限群 (Bieberbach 定理). (令 W 为由 \(W\) 生成的 \(\mathbf{R}^n\) 的向量子空间. 紧群 \(\Gamma \cap \mathbf{R}^n\)\(p(\Gamma)\) 只有有限个连通分支. 若 \(W = \{0\}\), 则由 (a), \(\Gamma\) 含有有限指数的交换正规子群 \((a)\)\(\Gamma_1\); 此时 \(\Gamma_1\) 也是晶体学群, 这与 (c) 矛盾. 因而 \((c)\)\(W \neq \{0\}\). 证明 \(p(\Gamma)\) 保持 W 不变且 \(W\)\(p(\Gamma)|W\) 是有限的; 否则存在  中的序列 \((S_m)\), 若 \(\Gamma\)\(S_m = (U_m, x_m)\), 则 \(U_m\) 彼此不同且趋于 I; 然后构造 \(I\)\((I, a_j)S_m(I, a_j)^{-1s_m^1}\), 其中 \((a_j)\) 是 \(\Gamma \cap \mathbf{R}^n\) 的一个基, 这将与 \(\Gamma\) 离散的假设矛盾. 最后, 为了看出 \(W = \mathbf{R}^n\), 证明否则 \(\Gamma\) 在 \(\mathbf{R}^n/W\) 上的作用将是一个不含非零平移的晶体学群的作用.)\(\neq 0\)
\end{enumerate}

